\documentclass[11pt]{article}
\usepackage[margin=2.6cm]{geometry}
\usepackage{mathpazo}
\usepackage{microtype}
\usepackage{booktabs}
\usepackage{graphicx}
\usepackage{float}
\usepackage{xcolor}
\usepackage{titlesec}
\usepackage{caption}
\usepackage{listings}
\usepackage[hidelinks]{hyperref}
\emergencystretch=3em

\definecolor{accent}{HTML}{1F4E79}
\definecolor{good}{HTML}{1A7F37}
\definecolor{bad}{HTML}{B42318}
\definecolor{warm}{HTML}{C9720B}
\titleformat{\section}{\Large\bfseries\color{accent}}{\thesection}{0.8em}{}
\titleformat{\subsection}{\large\bfseries\color{accent!80!black}}{\thesubsection}{0.8em}{}
\captionsetup{font=small,labelfont={bf,color=accent}}
\newcommand{\rA}{\textcolor{accent}{replica~101}}
\newcommand{\rB}{\textcolor{warm}{replica~202}}
\newcommand{\rC}{\textcolor{good}{replica~303}}

\lstdefinestyle{py}{
  basicstyle=\ttfamily\footnotesize,
  backgroundcolor=\color{gray!8},
  frame=single, rulecolor=\color{gray!40}, framerule=0.4pt,
  breaklines=true, columns=fullflexible, keepspaces=true,
  showstringspaces=false, aboveskip=8pt, belowskip=8pt}
\lstset{style=py}

\newsavebox{\eliBox}
\newenvironment{eli5}
  {\par\medskip\noindent\begin{lrbox}{\eliBox}\begin{minipage}{\dimexpr\linewidth-2\fboxsep-2\fboxrule}\small}
  {\end{minipage}\end{lrbox}\noindent\fcolorbox{accent!40}{accent!4}{\usebox{\eliBox}}\par\medskip}

% generated by make_numbers.py — do not edit
\newcommand{\rAModel}{google/gemini-3.5-flash}
\newcommand{\rAPrograms}{97}
\newcommand{\rAIncomplete}{4}
\newcommand{\rAValid}{93}
\newcommand{\rAInvalid}{4}
\newcommand{\rAValidPct}{95.9\%}
\newcommand{\rAWall}{20.5}
\newcommand{\rAWallWork}{19.1}
\newcommand{\rASeedCV}{0.848628}
\newcommand{\rASeedTest}{0.858852}
\newcommand{\rASeedRmse}{0.4290}
\newcommand{\rABestCV}{0.872291}
\newcommand{\rABestTest}{0.881128}
\newcommand{\rABestRmse}{0.3937}
\newcommand{\rABestId}{d36e1e62}
\newcommand{\rABestGen}{4}
\newcommand{\rABestNodes}{4}
\newcommand{\rABestDepth}{1}
\newcommand{\rABestFeats}{7}
\newcommand{\rADeltaCV}{+0.023662}
\newcommand{\rADeltaTest}{+0.022277}
\newcommand{\rADeltaRmse}{-0.0353}
\newcommand{\rAChildMean}{0.859426}
\newcommand{\rAChildStd}{0.004990}
\newcommand{\rANodesMean}{3.53}
\newcommand{\rANodesMax}{7}
\newcommand{\rADepthMean}{1.17}
\newcommand{\rADepthMax}{2}
\newcommand{\rADeepPct}{17.2\%}
\newcommand{\rACalls}{213}
\newcommand{\rACallsFailed}{2}
\newcommand{\rATokIn}{1,409,330}
\newcommand{\rATokOut}{1,359,318}
\newcommand{\rATokReason}{1,199,740}
\newcommand{\rATokReasonMean}{5,686}
\newcommand{\rATokOutMean}{6,442}
\newcommand{\rATokOutMax}{12,995}
\newcommand{\rALatMed}{40.8}
\newcommand{\rALatPfive}{58.4}
\newcommand{\rASpend}{14.35}
\newcommand{\rAMalformed}{0}
\newcommand{\rAInvalidTop}{own-target leakage: fit row 8 columns ['…'] change when only its own y}
\newcommand{\rANbrBuilt}{25}
\newcommand{\rANbrValid}{22}
\newcommand{\rANbrSup}{25}
\newcommand{\rANbrMention}{39}
\newcommand{\rANbrMentionOnly}{14}
\newcommand{\rANbrAgg}{38}
\newcommand{\rANbrFirst}{48}
\newcommand{\rANbrBest}{0.872291}
\newcommand{\rANbrChampion}{yes}
\newcommand{\rANbrGenFirst}{3}
\newcommand{\rAPanelN}{11}
\newcommand{\rAPanelNbr}{10}
\newcommand{\rAPanelRho}{+0.50}
\newcommand{\rAPanelTestMean}{0.881047}
\newcommand{\rAPanelTestMin}{0.875589}
\newcommand{\rAPanelTopTest}{0.883571}
\newcommand{\rAPanelTopTestId}{aab6fe26}
\newcommand{\rAPanelTopTestCVRank}{8}
\newcommand{\rAPanelRmseMean}{0.3938}
\newcommand{\rAPanelRmseWorst}{0.4028}
\newcommand{\rAPanelTopRmse}{0.3896}
\newcommand{\rAPanelTopRmseId}{aab6fe26}
\newcommand{\rAPanelTopRmseCVRank}{8}
\newcommand{\rANoNbrCV}{0.862380}
\newcommand{\rANoNbrTest}{0.875589}
\newcommand{\rANoNbrRmse}{0.4028}
\newcommand{\rANoNbrGen}{3}
\newcommand{\rANbrGapCV}{+0.009910}
\newcommand{\rANbrGapTest}{+0.005540}
\newcommand{\rANbrGapRmse}{-0.0091}
\newcommand{\rBModel}{google/gemini-3.5-flash}
\newcommand{\rBPrograms}{99}
\newcommand{\rBIncomplete}{2}
\newcommand{\rBValid}{96}
\newcommand{\rBInvalid}{3}
\newcommand{\rBValidPct}{97.0\%}
\newcommand{\rBWall}{19.0}
\newcommand{\rBWallWork}{17.7}
\newcommand{\rBSeedCV}{0.848628}
\newcommand{\rBSeedTest}{0.858852}
\newcommand{\rBSeedRmse}{0.4290}
\newcommand{\rBBestCV}{0.868428}
\newcommand{\rBBestTest}{0.878005}
\newcommand{\rBBestRmse}{0.3988}
\newcommand{\rBBestId}{e4c7cd13}
\newcommand{\rBBestGen}{6}
\newcommand{\rBBestNodes}{6}
\newcommand{\rBBestDepth}{1}
\newcommand{\rBBestFeats}{9}
\newcommand{\rBDeltaCV}{+0.019799}
\newcommand{\rBDeltaTest}{+0.019153}
\newcommand{\rBDeltaRmse}{-0.0302}
\newcommand{\rBChildMean}{0.859081}
\newcommand{\rBChildStd}{0.003983}
\newcommand{\rBNodesMean}{3.52}
\newcommand{\rBNodesMax}{7}
\newcommand{\rBDepthMean}{1.04}
\newcommand{\rBDepthMax}{2}
\newcommand{\rBDeepPct}{4.2\%}
\newcommand{\rBCalls}{197}
\newcommand{\rBCallsFailed}{2}
\newcommand{\rBTokIn}{1,326,207}
\newcommand{\rBTokOut}{1,132,899}
\newcommand{\rBTokReason}{997,158}
\newcommand{\rBTokReasonMean}{5,114}
\newcommand{\rBTokOutMean}{5,810}
\newcommand{\rBTokOutMax}{11,794}
\newcommand{\rBLatMed}{36.8}
\newcommand{\rBLatPfive}{56.8}
\newcommand{\rBSpend}{12.19}
\newcommand{\rBMalformed}{0}
\newcommand{\rBInvalidTop}{node spatial_clustering: output '…' declares numerical but has non-num}
\newcommand{\rBNbrBuilt}{14}
\newcommand{\rBNbrValid}{13}
\newcommand{\rBNbrSup}{14}
\newcommand{\rBNbrMention}{25}
\newcommand{\rBNbrMentionOnly}{11}
\newcommand{\rBNbrAgg}{42}
\newcommand{\rBNbrFirst}{63}
\newcommand{\rBNbrBest}{0.868428}
\newcommand{\rBNbrChampion}{yes}
\newcommand{\rBNbrGenFirst}{5}
\newcommand{\rBPanelN}{11}
\newcommand{\rBPanelNbr}{10}
\newcommand{\rBPanelRho}{-0.51}
\newcommand{\rBPanelTestMean}{0.878257}
\newcommand{\rBPanelTestMin}{0.875344}
\newcommand{\rBPanelTopTest}{0.880057}
\newcommand{\rBPanelTopTestId}{7a7334b5}
\newcommand{\rBPanelTopTestCVRank}{9}
\newcommand{\rBPanelRmseMean}{0.3984}
\newcommand{\rBPanelRmseWorst}{0.4032}
\newcommand{\rBPanelTopRmse}{0.3955}
\newcommand{\rBPanelTopRmseId}{7a7334b5}
\newcommand{\rBPanelTopRmseCVRank}{9}
\newcommand{\rBNoNbrCV}{0.864424}
\newcommand{\rBNoNbrTest}{0.879636}
\newcommand{\rBNoNbrRmse}{0.3962}
\newcommand{\rBNoNbrGen}{5}
\newcommand{\rBNbrGapCV}{+0.004004}
\newcommand{\rBNbrGapTest}{-0.001631}
\newcommand{\rBNbrGapRmse}{+0.0027}
\newcommand{\rCModel}{google/gemini-3.5-flash}
\newcommand{\rCPrograms}{100}
\newcommand{\rCIncomplete}{1}
\newcommand{\rCValid}{99}
\newcommand{\rCInvalid}{1}
\newcommand{\rCValidPct}{99.0\%}
\newcommand{\rCWall}{21.2}
\newcommand{\rCWallWork}{19.3}
\newcommand{\rCSeedCV}{0.848628}
\newcommand{\rCSeedTest}{0.858852}
\newcommand{\rCSeedRmse}{0.4290}
\newcommand{\rCBestCV}{0.873296}
\newcommand{\rCBestTest}{0.883771}
\newcommand{\rCBestRmse}{0.3893}
\newcommand{\rCBestId}{b4b25762}
\newcommand{\rCBestGen}{6}
\newcommand{\rCBestNodes}{6}
\newcommand{\rCBestDepth}{2}
\newcommand{\rCBestFeats}{12}
\newcommand{\rCDeltaCV}{+0.024667}
\newcommand{\rCDeltaTest}{+0.024920}
\newcommand{\rCDeltaRmse}{-0.0397}
\newcommand{\rCChildMean}{0.858588}
\newcommand{\rCChildStd}{0.004994}
\newcommand{\rCNodesMean}{3.54}
\newcommand{\rCNodesMax}{6}
\newcommand{\rCDepthMean}{1.09}
\newcommand{\rCDepthMax}{2}
\newcommand{\rCDeepPct}{9.1\%}
\newcommand{\rCCalls}{202}
\newcommand{\rCCallsFailed}{2}
\newcommand{\rCTokIn}{1,358,105}
\newcommand{\rCTokOut}{1,303,789}
\newcommand{\rCTokReason}{1,145,533}
\newcommand{\rCTokReasonMean}{5,728}
\newcommand{\rCTokOutMean}{6,519}
\newcommand{\rCTokOutMax}{13,165}
\newcommand{\rCLatMed}{40.9}
\newcommand{\rCLatPfive}{60.6}
\newcommand{\rCSpend}{13.77}
\newcommand{\rCMalformed}{0}
\newcommand{\rCInvalidTop}{target inverse: name '…' is not defined}
\newcommand{\rCNbrBuilt}{27}
\newcommand{\rCNbrValid}{26}
\newcommand{\rCNbrSup}{27}
\newcommand{\rCNbrMention}{36}
\newcommand{\rCNbrMentionOnly}{9}
\newcommand{\rCNbrAgg}{59}
\newcommand{\rCNbrFirst}{26}
\newcommand{\rCNbrBest}{0.873296}
\newcommand{\rCNbrChampion}{yes}
\newcommand{\rCNbrGenFirst}{2}
\newcommand{\rCPanelN}{11}
\newcommand{\rCPanelNbr}{10}
\newcommand{\rCPanelRho}{+0.65}
\newcommand{\rCPanelTestMean}{0.879406}
\newcommand{\rCPanelTestMin}{0.877397}
\newcommand{\rCPanelTopTest}{0.884570}
\newcommand{\rCPanelTopTestId}{35bf290d}
\newcommand{\rCPanelTopTestCVRank}{3}
\newcommand{\rCPanelRmseMean}{0.3965}
\newcommand{\rCPanelRmseWorst}{0.3998}
\newcommand{\rCPanelTopRmse}{0.3880}
\newcommand{\rCPanelTopRmseId}{35bf290d}
\newcommand{\rCPanelTopRmseCVRank}{3}
\newcommand{\rCNoNbrCV}{0.862100}
\newcommand{\rCNoNbrTest}{0.878315}
\newcommand{\rCNoNbrRmse}{0.3983}
\newcommand{\rCNoNbrGen}{4}
\newcommand{\rCNbrGapCV}{+0.011195}
\newcommand{\rCNbrGapTest}{+0.005456}
\newcommand{\rCNbrGapRmse}{-0.0090}
\newcommand{\allPanelScored}{33}
\newcommand{\allPanelNbr}{30}
\newcommand{\allPanelTopTest}{0.884570}
\newcommand{\allPanelTestMin}{0.875344}
\newcommand{\allPanelRho}{+0.66}
\newcommand{\allPanelBestRmse}{0.3880}
\newcommand{\allPanelWorstRmse}{0.4032}
\newcommand{\allPanelRmseRho}{+0.66}
\newcommand{\allNbrGapTestMean}{+0.003122}
\newcommand{\allNbrGapTestWins}{2}
\newcommand{\allNbrGapRmseMean}{-0.0051}
\newcommand{\allNbrGapRmseWins}{2}
\newcommand{\allReplicas}{3}
\newcommand{\allPrograms}{296}
\newcommand{\allValid}{288}
\newcommand{\allValidPct}{97.3\%}
\newcommand{\allNbrBuilt}{66}
\newcommand{\allNbrSup}{66}
\newcommand{\allNbrMention}{100}
\newcommand{\allNbrMentionOnly}{34}
\newcommand{\allNbrAgg}{139}
\newcommand{\allNbrReplicas}{3}
\newcommand{\allChampReplicas}{3}
\newcommand{\allBestCV}{0.873296}
\newcommand{\allBestTest}{0.883771}
\newcommand{\allBestRmse}{0.3893}
\newcommand{\allSpend}{40.30}
\newcommand{\allCalls}{606}
\newcommand{\priorPrograms}{200}
\newcommand{\priorNbrBuilt}{0}
\newcommand{\priorNbrMention}{1}
\newcommand{\priorNbrAgg}{13}

% generated by make_memory_report_assets.py — do not edit
\newcommand{\memProgramsStored}{101}
\newcommand{\memProgramsCompleted}{99}
\newcommand{\memProgramsValid}{94}
\newcommand{\memProgramsInvalid}{5}
\newcommand{\memProgramsIncomplete}{2}
\newcommand{\memWallMinutes}{140.2}
\newcommand{\memSeedCV}{0.850820}
\newcommand{\memBestCV}{0.881017}
\newcommand{\memBestCVStd}{0.004434}
\newcommand{\memDeltaCV}{+0.030196}
\newcommand{\memBestId}{777c20fb}
\newcommand{\memBestGeneration}{8}
\newcommand{\memBestNodes}{5}
\newcommand{\memBestDepth}{2}
\newcommand{\memBestFeatures}{33}
\newcommand{\memSeedTest}{0.858642}
\newcommand{\memSeedTestRmse}{0.429330}
\newcommand{\memBestTest}{0.891148}
\newcommand{\memBestTestRmse}{0.376746}
\newcommand{\memDeltaTest}{+0.032506}
\newcommand{\memDeltaTestRmse}{-0.052583}
\newcommand{\memExplicitBestId}{c1fbe7ff}
\newcommand{\memExplicitBestCV}{0.880501}
\newcommand{\memExplicitBestTest}{0.892790}
\newcommand{\memExplicitBestRmse}{0.373895}
\newcommand{\memNoMemoryBestLow}{0.868428}
\newcommand{\memNoMemoryBestHigh}{0.873296}
\newcommand{\memNoMemoryGainLow}{+0.019799}
\newcommand{\memNoMemoryGainHigh}{+0.024667}
\newcommand{\memCards}{52}
\newcommand{\memInsightCards}{41}
\newcommand{\memProgramCards}{11}
\newcommand{\memCardAdds}{52}
\newcommand{\memCardUpdates}{32}
\newcommand{\memCardRetired}{0}
\newcommand{\memDecisions}{109}
\newcommand{\memRandomized}{94}
\newcommand{\memDelivered}{70}
\newcommand{\memControls}{24}
\newcommand{\memExplicitUses}{36}
\newcommand{\memTerminals}{107}
\newcommand{\memEvidenceFinal}{92}
\newcommand{\memPendingFinal}{2}
\newcommand{\memEvidenceOnline}{89}
\newcommand{\memAllDeliveredGain}{+0.001850}
\newcommand{\memAllControlGain}{+0.002047}
\newcommand{\memLowDeliveredGain}{+0.004931}
\newcommand{\memLowControlGain}{+0.003890}
\newcommand{\memHighDeliveredGain}{-0.000053}
\newcommand{\memHighControlGain}{-0.000064}
\newcommand{\memSemanticSchemas}{1}
\newcommand{\memDynamicSchemas}{19}
\newcommand{\memRepeatedParents}{27}
\newcommand{\memMovedParents}{20}
\newcommand{\memCellErrors}{0}
\newcommand{\memCoordinateErrors}{0}
\newcommand{\memSemanticSpread}{0.0}
\newcommand{\memUnoccupied}{1}
\newcommand{\memSeMeasured}{93}
\newcommand{\memSeNonzero}{93}
\newcommand{\memFoundingSeZero}{69}
\newcommand{\memRetirementSweeps}{8}
\newcommand{\memBoundaryMass}{0.739}
\newcommand{\memBoundaryAlpha}{0.1}
\newcommand{\memGeminiCalls}{296}
\newcommand{\memQwenCalls}{383}
\newcommand{\memGeminiTokens}{4,970,543}
\newcommand{\memQwenTokens}{2,581,146}


\title{\color{accent}\bfseries The Capability Envelope and Memory V2 Follow-up\\
\large Deleting death classes from \texttt{dag\_tab}'s mutation surface:\\
can the search reach a supervised nearest-neighbour feature it could never reach before?\\
\large (historical no-memory replicas plus one current-main memory run, california)}
\author{GigaEvo}
\date{2026-07-22}

\begin{document}
\maketitle

\begin{abstract}
\noindent
An audit of the 2026-07-21 \texttt{dag\_tab} shakedown found that across
\priorPrograms{} mutations the search had \textbf{never once} built a
nearest-neighbour feature --- while a single hand-written supervised kNN node scored
0.865573~CV~$R^2$, beating every genome the search had ever produced. The diagnosis was not
that the mutator lacked the idea but that the surface killed it: a stack of static
AST rules, each a conservative approximation of Python semantics, rejected valid
neighbour code for spellings rather than for behaviour. This report measures what
happened after those rules were deleted and correctness moved to execution, with the
own-target leakage probe left as the only hard gate. Three replicas of 100 mutants each
were run on the identical recipe. \textbf{The capability appeared immediately and
everywhere.} Against \priorNbrBuilt{}~/~\priorPrograms{} before, the search built
\allNbrBuilt{} neighbour features across \allPrograms{} children, \emph{all
\allNbrSup{} of them supervised} --- and the champion of \textbf{all
\allChampReplicas{} of \allReplicas{} replicas} is a genome whose scoring path contains one.
The first valid one arrives at generation~\rCNbrGenFirst{} in \rC{}. Every replica also beat
the standing bar: best CV $R^2$ \rABestCV{}~/~\rBBestCV{}~/~\rCBestCV{} against the
2026-07-16 reference of 0.864554 and the hand-written oracle's 0.865573, and the climbs
transfer to the untouched test split: champion RMSE is
\rABestRmse{}~/~\rBBestRmse{}~/~\rCBestRmse{} against the seed's \rASeedRmse{}; the best
panel genome reaches \allPanelBestRmse{}. Two
findings qualify the headline. First, the search did not merely call a library: the winning
node in \rC{} implements \emph{leave-one-out} target encoding --- it requests $K{+}1$
neighbours and masks each fit row's own index --- which is the construction the leakage
probe demands, arrived at unprompted. Second, a residual death class remains and is now
measurable: \allNbrMentionOnly{} children \emph{proposed} a neighbour feature in prose and
shipped something else. Validity was \allValidPct{} over \allPrograms{} children, with
\textbf{zero} malformed or schema-invalid documents anywhere. Because selection ran against
CV, the obtained models were re-evaluated on the untouched test split as a panel rather than
as a single champion (\S\ref{sec:heldout}): the CV-top ten of \emph{every} replica is
\allPanelNbr{}/\allPanelNbr{} neighbour genomes, all \allPanelScored{} re-scored genomes beat
the seed out of sample, and the champion beats its own replica's best neighbour-free genome
in \allNbrGapRmseWins{} of \allReplicas{} replicas (mean RMSE gap
\allNbrGapRmseMean{}) --- it loses
that head-to-head in \rB{}, which we report rather than bury.

\smallskip\noindent
A follow-up run on current \texttt{main} exercises Memory V2 end to end with Gemini
3.5 Flash as mutator and Qwen Instruct as the memory model. Its champion reaches
\memBestCV{}~CV $R^2$ and \memBestTestRmse{} held-out RMSE; the best child that explicitly
cites a card reaches RMSE \memExplicitBestRmse{}. The historical no-memory trajectories are
shown as a reference envelope, not a causal ablation: the seed graph and evaluator revision
changed. The randomized ledger itself does not show a global delivery advantage
(\memAllDeliveredGain{} mean gain when delivered versus \memAllControlGain{} when
withheld), but it does show useful card-driven kNN mutations on weaker parents, stable
dynamic-MAP context, and nonzero evaluator uncertainty on every valid terminal.
\end{abstract}

\begin{eli5}
\textbf{In plain language.} The system invents new columns for a spreadsheet, and a
bolted-down model grades them. We knew one particular kind of column worked well --- ``look
at the ten geographically nearest houses and average what they sold for'' --- because a
human wrote it by hand and it beat everything the machine had found. But in two hundred
tries the machine had never written one. It turned out we had been rejecting that kind of
code for cosmetic reasons before ever running it: a pile of pattern-matching rules that
tried to guess what Python code would do, guessed wrong, and threw away good ideas. We
deleted the guessing and kept one real check --- the code must not peek at the answer for
the row it is scoring. In three fresh runs the machine wrote that kind of column sixty-six
times, won with it in all three runs, and worked out the anti-peeking trick by itself.
\end{eli5}

\section{The finding this run was built to test}

The prior shakedown was a success on its own terms: two arms, 100 mutants each, zero
malformed documents, a clean climb. The audit that followed asked a different question ---
not \emph{did it improve} but \emph{what could it never have found} --- and the answer was
specific. Re-running the same AST classifier used throughout this report over those two
run directories:

\begin{table}[H]\centering\small
\begin{tabular}{lrr}
\toprule
 & pre-refactor & this run \\
\midrule
children evaluated                       & \priorPrograms & \allPrograms \\
genomes computing a neighbour feature    & \textbf{\priorNbrBuilt} & \textbf{\allNbrBuilt} \\
of those, supervised (target-aware)      & 0 & \allNbrSup \\
genomes using the \texttt{aggregate} ABI at all & \priorNbrAgg & \allNbrAgg \\
genomes \emph{proposing} a neighbour feature in prose & \priorNbrMention & \allNbrMention \\
replicas whose champion contains one     & 0 & \textbf{\allChampReplicas~/~\allReplicas} \\
\bottomrule
\end{tabular}
\caption{The capability gap, measured rather than remembered. Both columns are classified by
the same detector (\S\ref{sec:detector}) over the stored genomes. The \texttt{aggregate}
row matters because only the aggregate ABI receives \texttt{y\_fit}: a supervised neighbour
feature is impossible without it, and before the refactor the search was barely using it.}
\end{table}

\noindent
The hand-written control that motivated all of this scores 0.865573~CV~$R^2$ as a single
node bolted onto the seed --- above the best genome any \texttt{dag\_tab} run had produced
at that point (0.864554). The idea was worth more than everything the search had found, and
the search had never tried it.

\subsection{What was actually blocking it}

The blocking rules were all instances of one mistake: \emph{a static approximation of Python
semantics is wrong in both directions, and one direction is always a death class}. A rule
that decides from the syntax tree whether a name still refers to the ABI frame, or whether a
read is undeclared, must be conservative somewhere; wherever it was conservative in the
rejecting direction, a valid runnable node died for a spelling. The refactor's rule was to
delete such a check or make it strictly conservative --- never to make it cleverer --- and
to let execution, which has the actual objects, be authoritative. Crucially, \textbf{nothing
about target leakage was relaxed}: \texttt{assert\_split\_invariant} still perturbs every
fit row's own target and requires that row's generated columns not move. The envelope was
widened; the correctness gate was not.

\section{Setup}

All three replicas ran from the review-fix worktree
(\texttt{fix/dag-tab-306-review-fixes-v2} at \texttt{bfa9d42c}), verified at launch by
resolving \texttt{gigaevo.\_\_file\_\_} and \texttt{problems/dag\_tab/validate.py} inside
that worktree. The recipe mirrors arm~B of the prior shakedown verbatim, so the only
intended difference from the pre-refactor column above is the mutation surface itself.

\begin{lstlisting}
python -u run.py problem.name=dag_tab program_format=json_document pipeline=guided \
  memory=none mutation=structured_diff_dag_tab \
  mutation_operator.allowed_changes.max_nodes=10 \
  llm=gemini35_flash llm.models.0.max_tokens=32768 \
  +llm.structured_output_method=function_calling \
  num_parents=1 max_mutants=100 algorithm=tabular/2d_local_ood
\end{lstlisting}

\noindent
\texttt{memory=none} and \texttt{num\_parents=1}: no memory apparatus, no crossover, so
every result here is attributable to the mutation operator alone. Fitness is 3-fold CV,
the tabular default --- the same protocol the 0.848628 seed and the 0.864554 bar were
measured under. The replicas differ only in LLM sampling; \texttt{gigaevo} has no global
RNG seed knob, so ``replica'' means an independent draw, not a seeded variant.

\subsection{How a neighbour feature is detected}
\label{sec:detector}

Counting the string \texttt{knn} would count a variable name as a capability. The detector
parses each node's body and classifies what the code \emph{computes}: \textbf{library} if it
imports or calls a neighbour primitive (\texttt{sklearn.neighbors}, \texttt{scipy.spatial},
pairwise distances); \textbf{handrolled} if it computes a distance \emph{and} ranks or
selects on it (\texttt{argsort}, \texttt{argpartition}, \texttt{searchsorted}); and
\textbf{none} otherwise, whatever the identifiers are called. It ships with four controls
covering both ways it could lie --- a variable named \texttt{knn} that computes a ratio, and
an \texttt{argsort} with no distance, both of which must classify as \texttt{none}. A node is
\emph{supervised} when it is an \texttt{aggregate} node, since only that ABI receives
\texttt{y\_fit}. No separate reachability test is needed: \texttt{FeatureGraph.estimator\_columns}
unions every node's \texttt{output\_cols} into the scored frame, so presence in a genome
\emph{is} presence on the scoring path.

\section{Did the search reach the capability?}

\begin{figure}[H]\centering
\includegraphics[width=\linewidth]{neighbour_funnel.png}
\caption{Left: how far along the funnel each replica got, from proposing a neighbour feature
in prose to shipping one that survives the leakage probe. The funnel is the point --- a run
that never mentions neighbours failed at retrieval, one that mentions and never builds failed
at translation, one that builds and is rejected failed at the probe, and collapsing these
into a single bar would hide which death class is still live. Right: the fitness of
neighbour genomes (rings) against everything else (faint dots), by generation.}
\end{figure}

\begin{table}[H]\centering\small
\begin{tabular}{lrrr}
\toprule
 & \rA & \rB & \rC \\
\midrule
children evaluated                    & \rAPrograms & \rBPrograms & \rCPrograms \\
proposes a neighbour feature in prose & \rANbrMention & \rBNbrMention & \rCNbrMention \\
uses the \texttt{aggregate} ABI       & \rANbrAgg & \rBNbrAgg & \rCNbrAgg \\
\textbf{computes a neighbour feature} & \textbf{\rANbrBuilt} & \textbf{\rBNbrBuilt} & \textbf{\rCNbrBuilt} \\
\quad of those, supervised            & \rANbrSup & \rBNbrSup & \rCNbrSup \\
\quad of those, survive the probe     & \rANbrValid & \rBNbrValid & \rCNbrValid \\
proposes but does not build           & \textcolor{bad}{\rANbrMentionOnly} & \textcolor{bad}{\rBNbrMentionOnly} & \textcolor{bad}{\rCNbrMentionOnly} \\
first valid one at generation         & \rANbrGenFirst & \rBNbrGenFirst & \rCNbrGenFirst \\
best neighbour genome (CV $R^2$)      & \rANbrBest & \rBNbrBest & \rCNbrBest \\
\textbf{champion contains one}        & \textbf{\rANbrChampion} & \textbf{\rBNbrChampion} & \textbf{\rCNbrChampion} \\
\bottomrule
\end{tabular}
\caption{Every replica reached the capability, and in every replica the capability won.
Notably \emph{all} \allNbrSup{} neighbour features built anywhere were supervised: having
been given a working aggregate ABI, the search went straight for the target rather than
settling for neighbour-averages of the inputs.}
\end{table}

\noindent
Two things in that table are worth separating from the headline.

\emph{The rate is not marginal.} This is not one lucky child in three hundred: between
\rBNbrBuilt{} and \rCNbrBuilt{} children per replica compute a neighbour feature, and
\rANbrValid{}, \rBNbrValid{} and \rCNbrValid{} of them respectively survive the leakage
probe. A capability that appears at this density in every replica is a property of the
surface, not of the draw.

\emph{One death class survives, and it is now visible.} \allNbrMentionOnly{} children
proposed a neighbour feature in their rationale and then shipped something else. That is no
longer the executor's doing --- these were never rejected --- so it is a failure of
translation inside the mutator, not of the envelope. It is the obvious next target, and it
was invisible before this run because there was nothing to compare the prose against.

\section{What the winning genome actually built}

The overall champion (\rC{}, \texttt{\rCBestId}, generation~\rCBestGen{}, CV~\rCBestCV{})
carries \emph{two} supervised neighbour nodes. This is its target encoder, verbatim:

\begin{lstlisting}
X_fit = df_fit.to_numpy();  X = df.to_numpy()
N = len(df_fit);  K = min(10, max(1, N - 1))

nn = NearestNeighbors(n_neighbors=K+1)      # K+1: room to drop the row itself
nn.fit(X_fit)
distances, indices = nn.kneighbors(X)

for i in range(len(df)):
    if i < N:                               # a fit row: leave itself out
        row_indices = indices[i]
        valid_indices = row_indices[row_indices != i][:K]
        pred[i] = y_fit[valid_indices].mean() if len(valid_indices) else y_mean
    else:                                   # held-out row: all K neighbours
        pred[i] = y_fit[indices[i][:K]].mean()

df['fe_spatial_knn_target'] = pred
\end{lstlisting}

\noindent
This is leave-one-out target encoding, and nothing in the prompt names it. The mutator asked
for one extra neighbour precisely so it could discard the row's own index on the fit side,
while using the full neighbourhood on the held-out side --- the exact asymmetry that makes
the feature honest and that \texttt{assert\_split\_invariant} tests for. Earlier in the same
replica, generation~2 produced a different valid solution to the same problem: a 5-fold
\texttt{KFold} out-of-fold encoder. The search found two distinct leak-free constructions
without being told either.

\begin{figure}[H]\centering
\includegraphics[width=\linewidth]{r303_lineage_graph.png}
\caption{\rC{} lineage, seed (left) to champion (right). Circles are nodes: blue
\texttt{rowwise}, orange \texttt{aggregate}. Arrows are \emph{consumed} edges --- a declared
dependency is drawn only where the child actually reads the parent's output columns. A red
ring marks a node added or modified at that step.}
\end{figure}

\section{Did it climb?}

\begin{figure}[H]\centering
\includegraphics[width=\linewidth]{ab_trajectory.png}
\caption{Best-so-far cross-validated $R^2$ against evaluation order (left) and wall clock
(right). The dashed red line is the 2026-07-16 reference on the pre-review code --- the bar
this refactor had to reach without regressing. All three replicas clear it, and clear the
hand-written kNN oracle (0.865573) as well.}
\end{figure}

\begin{table}[H]\centering\small
\begin{tabular}{lrrr}
\toprule
 & \rA & \rB & \rC \\
\midrule
children evaluated (incomplete)  & \rAPrograms~(\rAIncomplete) & \rBPrograms~(\rBIncomplete) & \rCPrograms~(\rCIncomplete) \\
valid / invalid                  & \rAValid~/~\rAInvalid & \rBValid~/~\rBInvalid & \rCValid~/~\rCInvalid \\
validity rate                    & \rAValidPct & \rBValidPct & \rCValidPct \\
malformed JSON / schema errors   & \textbf{\rAMalformed} & \textbf{\rBMalformed} & \textbf{\rCMalformed} \\
run duration (min)               & \rAWall & \rBWall & \rCWall \\
\midrule
seed CV $R^2$                    & \rASeedCV & \rBSeedCV & \rCSeedCV \\
best CV $R^2$                    & \rABestCV & \rBBestCV & \textbf{\rCBestCV} \\
$\Delta$ CV $R^2$                & \textcolor{good}{\rADeltaCV} & \textcolor{good}{\rBDeltaCV} & \textcolor{good}{\rCDeltaCV} \\
seed test RMSE                   & \rASeedRmse & \rBSeedRmse & \rCSeedRmse \\
best test RMSE                   & \rABestRmse & \rBBestRmse & \textbf{\rCBestRmse} \\
$\Delta$ test RMSE              & \textcolor{good}{\rADeltaRmse} & \textcolor{good}{\rBDeltaRmse} & \textcolor{good}{\rCDeltaRmse} \\
\midrule
valid-child CV mean $\pm$ sd     & \rAChildMean~$\pm$~\rAChildStd & \rBChildMean~$\pm$~\rBChildStd & \rCChildMean~$\pm$~\rCChildStd \\
champion (generation)            & \texttt{\rABestId}~(\rABestGen) & \texttt{\rBBestId}~(\rBBestGen) & \texttt{\rCBestId}~(\rCBestGen) \\
\quad nodes / depth / features   & \rABestNodes~/~\rABestDepth~/~\rABestFeats & \rBBestNodes~/~\rBBestDepth~/~\rBBestFeats & \rCBestNodes~/~\rCBestDepth~/~\rCBestFeats \\
nodes mean / max                 & \rANodesMean~/~\rANodesMax & \rBNodesMean~/~\rBNodesMax & \rCNodesMean~/~\rCNodesMax \\
consumed depth mean / max        & \rADepthMean~/~\rADepthMax & \rBDepthMean~/~\rBDepthMax & \rCDepthMean~/~\rCDepthMax \\
share at depth $\geq 2$          & \rADeepPct & \rBDeepPct & \rCDeepPct \\
\bottomrule
\end{tabular}
\caption{Test numbers come from \texttt{problems.dag\_tab.validate.score\_on\_test}, the
shipped delegation path, not a re-implementation in this report. The seed scores identically
in all three replicas, which is the sanity check that the delegated scorer is deterministic
across runs. Children still in the pipeline at teardown are excluded from every rate rather
than counted as failures.}
\end{table}

\noindent
\emph{Held-out error falls in every replica.} If the newly-permitted supervised features
were leaking the target into the CV objective, this is exactly where it would show: CV would
climb while test RMSE stayed flat or worsened. Instead it falls from \rASeedRmse{} to
\rABestRmse{}~/~\rBBestRmse{}~/~\rCBestRmse{}. Combined with the probe rejecting genuine
leakage attempts (\S\ref{sec:rejections}), that is the strongest evidence available here
that widening the envelope did not weaken the correctness gate.

\begin{figure}[H]\centering
\includegraphics[width=\linewidth]{ab_quality.png}
\caption{Validity counts, the fitness distribution over valid children, and the two
MAP-Elites behaviour axes the archive actually filled.}
\end{figure}

\section{Held-out evaluation of the obtained models}
\label{sec:heldout}

A champion's test RMSE is the error of the winner of a ${\sim}100$-way selection against
CV, and carries that optimism. So rather than scoring one genome per replica, the top ten
by CV were re-fitted and scored on the untouched test split --- plus, in each replica, the
best genome containing \emph{no} neighbour feature at all, which is the head-to-head the
capability claim actually rests on. Every number below comes from
\texttt{problems.dag\_tab.validate.score\_on\_test}, the shipped delegation path.

\begin{figure}[H]\centering
\includegraphics[width=\linewidth]{test_panel.png}
\caption{Left: CV $R^2$ (what selection optimised) against held-out RMSE for all
\allPanelScored{} re-scored genomes; filled circles contain a neighbour feature, hollow
squares do not, $\times$ is the seed. Right: per replica, the champion's test RMSE against
that replica's best neighbour-free genome. Lower RMSE is better.}
\end{figure}

\begin{table}[H]\centering\small
\begin{tabular}{lrrr}
\toprule
 & \rA & \rB & \rC \\
\midrule
genomes re-scored on test        & \rAPanelN & \rBPanelN & \rCPanelN \\
\quad of which contain a neighbour feature & \rAPanelNbr & \rBPanelNbr & \rCPanelNbr \\
test RMSE mean over the CV-top ten & \rAPanelRmseMean & \rBPanelRmseMean & \rCPanelRmseMean \\
test RMSE worst in panel         & \rAPanelRmseWorst & \rBPanelRmseWorst & \rCPanelRmseWorst \\
best test RMSE in panel          & \rAPanelTopRmse & \rBPanelTopRmse & \textbf{\rCPanelTopRmse} \\
\quad its CV rank                & \rAPanelTopRmseCVRank & \rBPanelTopRmseCVRank & \rCPanelTopRmseCVRank \\
CV-vs-test rank correlation ($\rho$, using $-\mathrm{RMSE}$) & \rAPanelRho & \textcolor{bad}{\rBPanelRho} & \rCPanelRho \\
\midrule
\multicolumn{4}{l}{\emph{best genome with NO neighbour feature}} \\
\quad CV $R^2$ (generation)      & \rANoNbrCV~(\rANoNbrGen) & \rBNoNbrCV~(\rBNoNbrGen) & \rCNoNbrCV~(\rCNoNbrGen) \\
\quad test RMSE                  & \rANoNbrRmse & \rBNoNbrRmse & \rCNoNbrRmse \\
\quad champion's CV lead over it & \textcolor{good}{\rANbrGapCV} & \textcolor{good}{\rBNbrGapCV} & \textcolor{good}{\rCNbrGapCV} \\
\quad champion's test RMSE gap (lower wins) & \textcolor{good}{\rANbrGapRmse} & \textcolor{bad}{\rBNbrGapRmse} & \textcolor{good}{\rCNbrGapRmse} \\
\bottomrule
\end{tabular}
\caption{Every genome in the panel is scored by the same delegated protocol as the seed. The
non-neighbour row is the appended best-neighbour-free genome, which is why the panel has
\rAPanelN{} rows rather than ten.}
\end{table}

\noindent
Three things fall out of this, and only two of them favour the headline.

\emph{The CV-top is entirely a neighbour-feature top.} Across all three replicas,
\textbf{\allPanelNbr{} of the \allPanelNbr{} CV-top-ten genomes contain a neighbour
feature} --- ten out of ten in every replica. The best neighbour-free genome in each run
sits at CV \rANoNbrCV{}~/~\rBNoNbrCV{}~/~\rCNoNbrCV{}, below the 0.864554 bar in two of the
three. On the objective the search actually optimised, the capability does not merely
appear: it dominates the top of the distribution.

\emph{The gains are real out of sample, but smaller than CV suggests.} Every one of the
\allPanelScored{} re-scored genomes beats the seed's \rASeedRmse{} RMSE on held-out data ---
the worst anywhere is \allPanelWorstRmse{}, and the best is \allPanelBestRmse{}. So this is
not a champion that got lucky; the whole CV-top of every run transfers. But the champion's
held-out advantage over the best neighbour-free genome is not uniform: its RMSE gaps are
\rANbrGapRmse{}, \rBNbrGapRmse{} and \rCNbrGapRmse{}.

\emph{And in one replica the head-to-head reverses.} In \rB{} the best neighbour-free
genome scores \rBNoNbrRmse{} RMSE against the champion's \rBBestRmse{} --- a
\rBNbrGapRmse{} loss for the neighbour genome, despite a \rBNbrGapCV{} CV lead. The
neighbour genome wins the head-to-head in \allNbrGapRmseWins{} of \allReplicas{} replicas,
with mean champion gap \allNbrGapRmseMean{}. Within the top band, CV differences of
${\sim}0.002$ simply do
not order test reliably: the rank correlation over the panel is \rAPanelRho{},
\rBPanelRho{} and \rCPanelRho{}, and in \rA{} the lowest test RMSE belongs to the genome
ranked \rAPanelTopRmseCVRank{}th on CV. The honest reading is that the neighbour capability
buys a solid gain over the seed and dominates the CV-top, and that \emph{among} strong
genomes the CV objective is too noisy at this resolution to rank them --- which is a
statement about the 3-fold objective, not about the capability.

\section{What the rejections were made of}
\label{sec:rejections}

\begin{table}[H]\centering\small
\begin{tabular}{lrrr}
\toprule
rejecting guard & replica 101 & replica 202 & replica 303 \\
\midrule
\textbf{malformed JSON / schema violation} & 0 & 0 & 0 \\
target transform/inverse: undefined or unbound name & 0 & 0 & 1 \\
target transform/inverse: wrong output shape & 0 & 0 & 0 \\
target transform/inverse: no round-trip & 0 & 0 & 0 \\
node: declared columns not produced & 0 & 0 & 0 \\
node: undeclared columns produced & 0 & 0 & 0 \\
node: declared numerical, produced non-numeric & 0 & 2 & 0 \\
node: removed pre-existing columns & 1 & 0 & 0 \\
own-target leakage (probe) & 2 & 1 & 0 \\
node body: Python error at execution & 1 & 0 & 0 \\
\midrule
total invalid & 4 & 3 & 1 \\
\bottomrule
\end{tabular}

\caption{Every invalid child, bucketed by the guard that rejected it. The first row is the
one the grammar owns and it is empty everywhere: across \allPrograms{} children, not one was
rejected for being unparseable or schema-invalid. Every rejection was semantic.}
\end{table}

\noindent
Total invalid across all three replicas is small enough to enumerate individually, and the
composition is the interesting part. \textbf{The leakage probe fired three times} --- on
children that attempted a supervised feature and got the leave-one-out construction wrong.
That is the gate doing precisely its job: the same widening that let \allNbrBuilt{} honest
neighbour features through also let dishonest ones be attempted, and they were caught at
execution rather than waved through. Two more were rejected for emitting a non-numeric
column from a clustering node, one for an unbound name in a target inverse, one for a node
body raising at runtime, and one for removing pre-existing columns.

\section{What it cost}

\begin{figure}[H]\centering
\includegraphics[width=\linewidth]{ab_cost.png}
\caption{Token budget per successful call, mutation-call latency, list-price spend, and wall
clock per replica. All three ran \texttt{gemini-3.5-flash} at \texttt{reasoning.effort:high}
under \texttt{service\_tier:flex}, so the latency distribution is the sum of thinking time
and flex queueing and should not be read as a pure thinking cost.}
\end{figure}

\begin{table}[H]\centering\small
\begin{tabular}{lrrr}
\toprule
 & \rA & \rB & \rC \\
\midrule
LLM calls (failed)       & \rACalls~(\rACallsFailed) & \rBCalls~(\rBCallsFailed) & \rCCalls~(\rCCallsFailed) \\
input tokens             & \rATokIn & \rBTokIn & \rCTokIn \\
output tokens            & \rATokOut & \rBTokOut & \rCTokOut \\
reasoning tokens         & \rATokReason & \rBTokReason & \rCTokReason \\
\quad per call           & \rATokReasonMean & \rBTokReasonMean & \rCTokReasonMean \\
output mean / max        & \rATokOutMean~/~\rATokOutMax & \rBTokOutMean~/~\rBTokOutMax & \rCTokOutMean~/~\rCTokOutMax \\
latency median / p95 (s) & \rALatMed~/~\rALatPfive & \rBLatMed~/~\rBLatPfive & \rCLatMed~/~\rCLatPfive \\
list-price spend (USD)   & \$\rASpend & \$\rBSpend & \$\rCSpend \\
\bottomrule
\end{tabular}
\caption{List-price spend at OpenRouter rates (\$1.50/\$9.00 per~M for
\texttt{gemini-3.5-flash}); thinking tokens bill at the completion rate.
\texttt{service\_tier:flex} discounts the actual invoice below these figures. All three
replicas together cost \$\allSpend{} at list price over \allCalls{} successful calls.}
\end{table}

\noindent
The 32768-token cap was never binding: the longest completion anywhere was
\rCTokOutMax{} tokens. Running three replicas concurrently on one machine cost roughly the
wall clock of one.

\section{Memory V2 follow-up on current main}
\label{sec:memoryv2}

This follow-up is a systems run and a source of hypotheses, not a matched memory A/B.
It used the neutral \texttt{dag\_tab\_seed}, the current aggregate-node evaluator and the
current structured-diff surface; the three no-memory trajectories above came from
\texttt{bfa9d42c}, with a one-node seed and an earlier evaluator. What is controlled is the
problem, fixed CatBoost scorer, 100-mutant budget, mutation model
(\texttt{google/gemini-3.5-flash}, high reasoning, flex tier), one-parent structured diff,
ten-node cap and MAP-Elites recipe. Memory selection, retrieval and authoring used
\texttt{Qwen3-235B-A22B-Instruct}. The run occupied \memWallMinutes{} minutes and made
\memGeminiCalls{} Gemini calls (\memGeminiTokens{} tokens) plus \memQwenCalls{} memory-model
calls (\memQwenTokens{} tokens). \texttt{memory.run\_seed=101} seeds assignment replay; it
does not seed Gemini sampling and does not make the historical directory named
\texttt{r101} a paired control.

\subsection{Convergence and untouched-test result}

\begin{figure}[H]\centering
\includegraphics[width=\linewidth]{memory_convergence.png}
\caption{Best-so-far CV $R^2$ for the three historical no-memory Gemini 3.5 Flash replicas
and the single current-main Memory V2 run. The right panel subtracts each run's own seed,
which removes the initial 0.0022 score shift but cannot remove evaluator and mutation-surface
changes. The red curve is therefore evidence that the complete memory-enabled stack climbs,
not an estimate of memory's causal effect.}
\end{figure}

\begin{table}[H]\centering\small
\begin{tabular}{lrrrr}
\toprule
run & seed CV $R^2$ & best CV $R^2$ & $\Delta$ CV & best test RMSE \\
\midrule
historical no-memory 101 & \rASeedCV & \rABestCV & \rADeltaCV & \rABestRmse \\
historical no-memory 202 & \rBSeedCV & \rBBestCV & \rBDeltaCV & \rBBestRmse \\
historical no-memory 303 & \rCSeedCV & \rCBestCV & \rCDeltaCV & \rCBestRmse \\
\textbf{Memory V2} & \memSeedCV & \textbf{\memBestCV} & \textbf{\memDeltaCV} & \textbf{\memBestTestRmse} \\
\bottomrule
\end{tabular}
\caption{The test entry is the CV champion re-fitted through the shipped
\texttt{score\_on\_test} path. Historical rows retain their original run results and are
not revision-matched controls.}
\end{table}

The memory run stored \memProgramsStored{} programs: \memProgramsCompleted{} completed,
\memProgramsValid{} valid, \memProgramsInvalid{} evaluated-invalid and
\memProgramsIncomplete{} still in flight when the hard cap fired. The generation
\memBestGeneration{} champion \texttt{\memBestId} has CV
\memBestCV{}~$\pm$~\memBestCVStd{} and test RMSE \memBestTestRmse{}, improvements of
\memDeltaCV{} CV $R^2$ and \memDeltaTestRmse{} RMSE over its neutral seed
(RMSE \memSeedTestRmse{}).
The best explicitly card-assisted genome, \texttt{\memExplicitBestId}, has slightly lower
CV (\memExplicitBestCV{}) but lower test RMSE (\memExplicitBestRmse{}). That reversal is
small but useful: 3-fold CV can choose the final
increment incorrectly even when the lineage as a whole improves strongly.

\subsection{What the memory system did}

\begin{figure}[H]\centering
\includegraphics[width=\linewidth]{memory_dynamics.png}
\caption{Top left: writer refreshes; the shaded band is pending evidence. Top right:
valid randomized proposals, with the evaluator-reported standard error rather than a
synthetic \texttt{se=0}. Bottom left: raw intention-to-treat means by parent-fitness band.
Bottom right: assignment, post-treatment citation and validity. Explicit citation is useful
for mechanism inspection, but only delivered-versus-withheld assignment is randomized.}
\end{figure}

The causal ledger contains \memDecisions{} decisions and \memTerminals{} terminals.
After the cold-start/empty decisions, \memRandomized{} card proposals were randomized:
\memDelivered{} delivered and \memControls{} withheld. The mutator explicitly grounded a
card in \memExplicitUses{} outcomes. The final writer snapshot released
\memEvidenceFinal{} eligible observations and retained \memPendingFinal{} administratively
cancelled mutations as pending, not invalid.

The global raw intention-to-treat result is deliberately unexciting: mean valid fitness
gain is \memAllDeliveredGain{} under delivery and \memAllControlGain{} under a withheld
card. This run does \emph{not} establish a global causal win for memory. There is a useful
interaction hypothesis, however. Below parent fitness 0.865 the corresponding means are
\memLowDeliveredGain{} and \memLowControlGain{}; at 0.875 or above both are effectively
flat (\memHighDeliveredGain{} and \memHighControlGain{}). That agrees with the observed
role of kNN cards: they frequently rescue weak spatial programs, while near the frontier
most additional mutations --- guided or not --- have little room and often regress. The
fitness bands are a post-hoc slice of one run, so this is a target for a matched replicate,
not a confirmatory subgroup result.

\begin{table}[H]\centering\small
\begin{tabular}{p{5.1cm}lrrrr}
\toprule
card mechanism & kind & proposed & control & cited & best cited $\Delta R^2$ \\
\midrule
multi-scale kNN target means and micro/macro ratio \newline \scriptsize{mem-45c42d56} & insight & 2 & 1 & 1 & +0.003444 \\
add local volatility, density, and covariate means \newline \scriptsize{mem-b491828e} & insight & 6 & 2 & 2 & +0.006621 \\
out-of-fold kNN target embedding \newline \scriptsize{mem-c27d9dc3} & insight & 4 & 1 & 3 & +0.005789 \\
K-means centroid-distance features \newline \scriptsize{mem-8adc6313} & insight & 5 & 0 & 2 & +0.004900 \\
normalized distance-to-economic-hub ratios \newline \scriptsize{mem-82d70d2d} & insight & 3 & 0 & 2 & +0.003673 \\
program exemplar: rotations, landmarks, clustering \newline \scriptsize{program-f4e6} & program & 6 & 3 & 2 & +0.002913 \\
\bottomrule
\end{tabular}

\caption{Representative active cards. ``Cited'' means the structured mutation explicitly
returned the delivered id; its gain is descriptive post-treatment evidence, not a
card-level randomized effect. Sparse arms are why no per-card causal contrast is printed.
The enriched-neighbourhood card produced two closely replicated gains of approximately
0.0066 and 0.0062, while the out-of-fold kNN and K-means cards each produced useful
low-fitness children.}
\end{table}

The bank ended at \memCards{} cards (\memInsightCards{} insights and
\memProgramCards{} program exemplars) after \memCardAdds{} additions and
\memCardUpdates{} equivalence/evidence updates. No card retired. Retirement was attempted
but disabled in \memRetirementSweeps{} refreshes because the immediate-reward residual-scale
posterior placed \memBoundaryMass{} mass at its configured boundary, above
$\alpha=\memBoundaryAlpha{}$. This is conservative --- no card was silently removed on a
bad fit --- but a 52-card bank after 100 mutations is too fast-growing for a long campaign.
The residual-scale bounds or retirement calibration should be tuned before treating this
bank size as healthy.

\subsection{Dynamic MAP cells and evaluation uncertainty}

The dynamic axes really did move. Across \memDecisions{} frozen decisions there are
\memDynamicSchemas{} live behavior-space hashes but only \memSemanticSchemas{} stable
semantic schema. Of \memRepeatedParents{} parents selected repeatedly,
\memMovedParents{} changed discrete cell at least once as the bounds expanded. A concrete
parent moved through cells $(13,9)\to(8,8)\to(7,6)\to(6,6)\to(6,5)$ without its raw behavior
changing.

That movement is safe because the posterior does \textbf{not} regress on the transient bin
id or live dynamic-normalized coordinate. It uses the fixed
\texttt{semantic\_normalized} coordinate; live bounds, bin and cell remain frozen audit and
archive-locality state. The repeated-parent semantic spread is \memSemanticSpread{}, with
\memCellErrors{} cell/coordinate alignment errors and \memCoordinateErrors{} out-of-bound
dynamic coordinates. Archive reindexing publishes a complete cell replacement, and the
memory snapshot copies elites and bounds under the same island lock. One decision records
\texttt{parent\_cell\_occupied=false}: the selected parent was no longer the incumbent by
snapshot time under steady-state concurrency, but its cell and bounds are internally
consistent and the occupancy bit is not a posterior feature.

Evaluator uncertainty also reaches the causal ledger. All \memSeMeasured{} valid terminal
measurements carry an SE and all \memSeNonzero{} are nonzero; the plotted bars come directly
from cross-validation fold SDs. The live process began before the final founding-card audit
fix, so \memFoundingSeZero{} historical \texttt{cards.json} founding events still show
\texttt{gain\_se=0}. That field did not train the Bayesian model: SQLite terminal
measurements are authoritative and contain the nonzero values. Commit \texttt{dc952fda}
now preserves the same uncertainty in newly authored card audit events as well.

The last online choice was fit with \memEvidenceOnline{} observations; the final writer
reached \memEvidenceFinal{} only after that choice. We therefore report ledger outcomes and
historical context-conditional predictions, not a purported ``final posterior'' copied
from the last decision. Replaying all \memEvidenceFinal{} rows at a fixed context is the
right follow-up if a final card-ranking posterior is needed.

\subsection{What the winning feature DAG built}

\begin{figure}[H]\centering
\includegraphics[width=\linewidth]{memory_champion_dag.png}
\caption{The champion's internal feature DAG. Four feature blocks consume raw columns
directly. Only \texttt{neighborhood\_ratios} consumes a generated block, making the graph
depth two. Every block is an estimator output, and the eight raw columns remain present.}
\end{figure}

\begin{table}[H]\centering\small
\begin{tabular}{p{3.8cm}p{3.0cm}rp{6.0cm}}
\toprule
node (kind) & dependency & outputs & role \\
\midrule
geo rotations (rowwise) & raw only & 6 & coast-aligned coordinate projections \\
geo distances (rowwise) & raw only & 6 & metro proximity and bounded regional gradients \\
demographic ratios (rowwise) & raw only & 5 & household scale and occupancy-normalized composition \\
spatial neighborhood (aggregate) & raw only & 11 & leave-one-out local target/covariate summaries at three scales \\
neighborhood ratios (rowwise) & spatial neighborhood & 5 & household-to-local and micro-to-macro contrasts \\
\bottomrule
\end{tabular}

\caption{All five champion nodes. Output counts sum to \memBestFeatures{} generated
features; no target transform or raw-column drop is used.}
\end{table}

The central feature is \texttt{spatial\_neighborhood}. It fits a nearest-neighbour index on
latitude/longitude, requests up to 65 points, removes a fitting row's own index, and emits
target means at $k=4,16,64$, target spread at $k=16$, three density proxies and four local
covariate means. The dependent node compares raw households with those local means and
forms a micro-to-macro target ratio. This is the high-value path. Rotated coordinates make
diagonal geography easier for axis-aligned trees; Haversine landmark distances encode
proximity to LA, SF, San Diego and San Jose; demographic ratios expose household scale and
occupancy. The latter blocks are plausible complementary signals, although fixed-city
landmarks are California-specific and several rotations are redundant in a linear sense.

The lineage explains the attribution. It climbs from \memSeedCV{} to 0.862772 with three
rowwise geographic/demographic blocks, adds supervised spatial kNN at generation~3, then
the ratio dependency and richer local statistics. Generation~6 explicitly cites a
distance-ratio card; generation~7 explicitly cites
\texttt{mem-45c42d56} to replace a single-scale target encoder with the multi-scale
construction, gaining +0.003444 to \memExplicitBestCV{}. The generation~8 champion was
offered \texttt{mem-9dfd62e3} but did not cite it; its +0.000516 demographic expansion came
from intra-lineage/program guidance. Memory therefore materially shaped the winning
lineage, but the final edge itself is correctly recorded as delivered-and-ignored, not
credited to a card.

\section{Threats to validity and known limits}

\begin{itemize}
\item \textbf{The Memory V2 comparison is historical and unmatched.} It is one current-main
run against three older no-memory runs with a different neutral seed, evaluator and mutation
surface. Normalizing by each seed does not repair that. A causal system-level memory effect
requires fresh no-memory replicas on the same commit and seed graph.
\item \textbf{Card citation is post-treatment.} The offered-versus-withheld contrast is
randomized; ``explicitly used'' is a mutator response to seeing the card. Per-card cited
gains and the frontier bands diagnose mechanisms but are not randomized subgroup estimates.
\item \textbf{No matched control on unmodified \texttt{main}.} The pre-refactor column comes
from the 2026-07-21 shakedown, which used the same recipe and the same model but a different
code revision, and was one run rather than three. The capability gap
(\priorNbrBuilt{}~$\to$~\allNbrBuilt{}) is far too large to be a sampling artefact, but the
\emph{fitness} comparison across that boundary mixes the refactor with run-to-run variance.
Three control replicas on unmodified \texttt{main} would settle it and have not been run.
\item \textbf{The leakage probe runs on 1024 rows; production fits ${\sim}14$k.} Own-target
invariance is verified on the first 1024 rows of \texttt{X\_train}. A dependence that only
manifests at larger fit sizes would not be caught. This is a known, disclosed residual of the
probe design, unchanged by this refactor.
\item \textbf{Test RMSE is reported once, at the end, but CV $R^2$ drove selection.} With
${\sim}100$ children per replica drawn against a CV objective, part of every CV gain is
selection noise. Held-out RMSE is the honest number and is the one quoted in the abstract.
\item \textbf{``Supervised'' is inferred from the ABI, not from dataflow.} A node is counted
as supervised when it is an \texttt{aggregate} node, which is the only kind that receives
\texttt{y\_fit}. An aggregate node that ignores its target would be over-counted. Every
neighbour node inspected by hand did use it, but the classifier does not prove it.
\item \textbf{California is one dataset, and a spatial one.} Nearest-neighbour features are
unusually well-suited to geographic data. That the envelope now \emph{admits} them is a
framework result; that they \emph{win} is a statement about california.
\item \textbf{The reviewer verdict on the refactor is \emph{would not launch}, and these
runs launched anyway.} The round-13 correctness review landed after the runs and reproduced
two blocking defects, both in the same component --- literal transcription of a node's frame
reads into its \texttt{input\_cols}. A \texttt{match}/\texttt{case} capture that rebinds the
frame name, and a temporary column that happens to share a raw column's name, each cause a
raw column to be added to \texttt{input\_cols}; the node is then rejected for ``removed
pre-existing columns''. \emph{Both defects point in the false-rejection direction}: they kill
valid runnable nodes, so they can only have narrowed the envelope measured here, never
inflated it --- the one such rejection in this campaign appears in the invalid table. The
capability counts are therefore a lower bound on what the refactor admits, but the refactor
is not shippable as reviewed, and any number here may move when the two defects are fixed.
\item \textbf{Held-out re-scoring re-fits the estimator.} The panel in
\S\ref{sec:heldout} re-runs each graph through \texttt{score\_on\_test}, which fits on train
and scores on test. CatBoost is seeded but not bitwise-deterministic across machines, so the
final digit of any test figure is not meaningful.
\end{itemize}

\section{What this run establishes}

\begin{enumerate}
\item The capability gap was caused by the mutation \emph{surface}, not by the mutator.
Nothing about the model, prompt, recipe or budget changed between
\priorNbrBuilt{}~/~\priorPrograms{} and \allNbrBuilt{}~/~\allPrograms{} neighbour features ---
only the removal of static rules that rejected valid code before it ran.
\item Widening the envelope did not weaken the correctness gate. The leakage probe fired on
three genuine attempts, held-out RMSE decreased in every replica, and the winning genome
independently implemented leave-one-out encoding to satisfy a probe it was never shown.
\item The capability translates into fitness, and survives held-out re-scoring. All
\allChampReplicas{} champions contain a supervised neighbour feature, the CV-top ten of every
replica is entirely neighbour genomes, all three replicas cleared both the 0.864554 standing
bar and the 0.865573 hand-written oracle, and every one of the \allPanelScored{} re-scored
genomes beats the seed on the untouched test split (best RMSE \allPanelBestRmse{}). The
neighbour-versus-no-neighbour head-to-head goes \allNbrGapRmseWins{}--1 on test with mean
RMSE gap \allNbrGapRmseMean{}, so the capability's advantage over strong neighbour-free genomes is
real but not yet separated from noise by three replicas.
\item A residual death class is now measured rather than suspected:
\allNbrMentionOnly{} children proposed a neighbour feature in prose and shipped something
else. That failure lives inside the mutator, not the envelope, and is the next thing to fix.
\item The current-main Memory V2 stack completed a full 100-mutant campaign with stable
dynamic-MAP context and evaluator uncertainty in every usable Bayesian outcome. Its winning
lineage uses cards to build the multi-scale nearest-neighbour path and reaches
\memBestTestRmse{} held-out RMSE (best explicitly assisted genome
\memExplicitBestRmse{}).
The randomized ledger does not support a global memory-win claim, and the non-retiring,
fast-growing card bank is the clearest operational issue exposed by the run.
\end{enumerate}

\end{document}
